\section{Introduction}
\label{sec:introduction}

Private compatibility can make formal incompatibility less technologically important without making formal harmonization less politically attractive. A firm may redesign a product, obtain an additional certification, or support an additional interface so that it can serve markets governed by another standard. Such a workaround narrows the effective technological difference between incompatible regimes. But it also changes the rents that governments obtain from maintaining an exclusive formal arrangement. The same private action can therefore substitute for harmonization technologically while changing---in either direction---the government's incentive for broader formal coordination.

This paper asks when that political effect is positive, when it is too small to change the government's ranking, and when it runs in the opposite direction. The answer is conditional. In the canonical three-country Cournot microfoundation, a member of a regional standards union can initially prefer the regional arrangement to international standardization because the regional standard creates an advantage in member markets, even though the member remains disadvantaged in the outsider market. If the outsider privately adopts the bloc standard, the member-market advantage is eroded. Whether the member then moves toward international standardization depends on how much rent remains after adaptation.

The central object is an equilibrium-derived residual-rent function. Let $d=\lambda c$, with $0\leq\lambda\leq1$, denote the marginal adaptation cost that remains after the outsider adopts the bloc standard and becomes fully compatible in member markets. The regional member's pre-adoption payoff difference relative to international standardization is
\begin{equation}
    W_M^{SU,N}-W^{IS}=\mathscr E-\mathscr D,
    \label{eq:intro-pre-gap}
\end{equation}
where $\mathscr E$ is the net member-market exclusion component and $\mathscr D$ is the reciprocal disadvantage in the outsider market. Re-solving the product-market equilibrium after incomplete adaptation gives
\begin{equation}
    W_M^{SU,O}(d)-W^{IS}
    =-\mathscr D+R(d,v),
    \label{eq:intro-post-gap}
\end{equation}
with
\begin{equation}
    R(d,v)
    =
    \frac{d\{(5-4v)d+2(1-4v)\}}
         {32(1-v)^2}.
    \label{eq:intro-residual-rent}
\end{equation}
On the canonical network domain, $R(d,v)$ is increasing in $d$. Conditional on $\mathscr E>\mathscr D>0$, the comparison is therefore exact: if $R(d,v)<\mathscr D$, private adaptation reverses the member government's ranking from the regional union to international standardization; if $\mathscr D<R(d,v)<\mathscr E$, the incentive for international standardization becomes stronger but the regional union remains preferred; and if $R(d,v)>\mathscr E$, private adaptation weakens the relative incentive for international standardization. The technological response does not have a sign-invariant political effect.

The result is not a knife-edge consequence of complete cost elimination. In the partial-adaptation model, outsider-only adoption and preference reversal coexist on a nonempty open set with $0<\lambda<1$. At the exact witness
\[
(c,v,\lambda,F)
=
\left(\frac1{10},\frac6{25},\frac12,\frac9{100}\right),
\]
one-half of the original marginal adaptation cost remains, yet the outsider strictly adopts the bloc standard, coalition members strictly do not reciprocate, separate national standards induce no private adoption, and a regional member strictly switches its ranking in favor of international standardization. The fixed-cost adoption conditions and the government-ranking conditions are thus satisfied for the same primitives rather than in separate examples.

This government-incentive result rests on a private-adoption continuation that is itself endogenous. A standard covers all markets in the formal coalition, so the return to one adoption can span several destinations. In the symmetric regional union, the outsider may obtain the value of supporting the bloc standard in two member markets while each member's reverse adoption affects only the outsider market. More generally, a selection-free outsider-only continuation is guaranteed when the outsider's worst-case adoption gain is positive and every member's best-case reverse-adoption gain is negative. In the canonical model these inequalities reduce to explicit fixed-cost thresholds. We use this adoption structure as supporting machinery rather than as a claim that one-way compatibility is itself novel.

The institutional application is deliberately narrower than the government-incentive result. Under the full-bypass endpoint $d=0$, and on the certified canonical parameter region, a high private-adoption cost leaves all three symmetric regional standards unions stable under strict blocking. For an intermediate fixed-cost range, outsider-only adoption makes all three countries strictly prefer international standardization to the continuation payoff under any regional union, and international standardization is uniquely stable. This stable-set comparison illustrates how a private technological response can alter the continuation payoffs supporting a formal coalition.

The high-cost stable set, however, is not robust in the same sense as the intermediate unique-international result. Under weak/Pareto blocking, symmetry-induced indifference allows one regional union to weakly block another, so the symmetric high-cost stable set is empty. A small market-size asymmetry breaks that indifference and selects the regional union formed by the two larger markets. By contrast, the intermediate unique-international result is supported by strict payoff differences and survives both blocking concepts under the certified small-asymmetry perturbation. Coalition stability is therefore an application of the payoff mechanism, not evidence that the particular high-cost stable-set correspondence is institution free.

The paper also reports where the political mechanism fails. First, at the zero-network boundary of the canonical model,
\[
W_M^{SU,N}-W^{IS}
=
\frac{c(13c-6)}{32}<0
\qquad (0<c<1/3),
\]
so the regional union is not initially preferred even though one-way private adoption can still occur. Second, a pre-specified alternative that grants singleton firms their own network benefit eliminates the initial regional advantage throughout the old canonical domain while preserving a one-way-adoption witness. Third, in a pre-specified differentiated-demand environment, the initial regional advantage and the positive member-market incentive term fail throughout the audited parameter box under both Cournot and Bertrand competition, even though outsider-only private adoption remains feasible. The portability failure therefore cannot be attributed to price competition alone. These negative results are part of the contribution because they separate the relatively portable adoption logic from the more specification-sensitive government-ranking effect.

The paper relates to established work on converters, interoperability, endogenous compatibility, costly multi-standard production, standards policy, and regional versus multilateral harmonization. Private workarounds to incompatibility are well known \citep{farrellSaloner1992,doganogluWright2006}, and one-way or endogenous compatibility can arise with strategic firm choice \citep{buccellaEtAl2023}. Fixed-cost economies of scope and dependent adoption decisions likewise have broad antecedents. On the policy side, standards unions, strategic recognition, and regional-versus-multilateral standards comparisons are already studied in international trade and standards policy \citep{gandalShy2001,takaradaEtAl2020,kawabataTakarada2021}. Recent work also emphasizes endogenous firm adoption of harmonized standards and the possibility of spontaneous regulatory harmonization \citep{schmidtSteingress2022,maggiMrazova2024}. The contribution here is therefore not any one ingredient. It is the explicit conditional mapping from a post-formation private adaptation decision to the continuation rents that enter a government's comparison of formal standards regimes, together with the success and failure boundaries of that mapping.

The analysis proceeds in three layers. Sections~\ref{sec:model} and \ref{sec:product-market} define the canonical model and collect the Cournot building blocks. Section~\ref{sec:private-compatibility} characterizes the private-adoption continuation and its selection-free region. Section~\ref{sec:selective-erosion} derives the residual-rent government-incentive result and the incomplete-adaptation open set. Section~\ref{sec:coalition-stability} applies the mechanism to formal coalition stability. Section~\ref{sec:secondary} then reports the network, demand, and institutional robustness boundaries. The discussion distinguishes the model's conditional mechanism from broader claims that the analysis does not establish.
